% QuorumGraph -- sprint submission report.
% Self-contained: all tables are inline; the only other file needed is references.bib.
% Build: pdflatex report && bibtex report && pdflatex report && pdflatex report  (or latexmk -pdf report.tex)
% Page budget: main text (Sections 1-6) must stay within 8 pages; references and appendices do not count.
% Every number below comes from docs/report/tables/numbers.log (generated by docs/report/numbers.py).
%
% Items marked \todo{...} render in red on purpose so they cannot ship by accident.
\documentclass[11pt]{article}
\usepackage[margin=1in]{geometry}
\usepackage[T1]{fontenc}
\usepackage{lmodern}
\usepackage{booktabs}
\usepackage{amsmath}
\usepackage[numbers]{natbib}
\usepackage{xcolor}
\usepackage[hidelinks]{hyperref}
\setlength{\emergencystretch}{3em}
\newcommand{\ci}[2]{\,{\scriptsize\color{gray}[#1,#2]}}
\newcommand{\code}[1]{\texttt{#1}}
\newcommand{\todo}[1]{{\color{red}\textbf{[TODO: #1]}}}

\title{QuorumGraph: Alignment That Survives Self-Improvement\\[4pt]
\large Evidence-based trust for self-modifying agents, tested by evolutionary simulation through a real trust service}
\author{Evan Zhang\\ \todo{affiliation}}
\date{}

\begin{document}
\maketitle

\begin{abstract}
An agent that improves itself produces a successor that is a new system: alignment evidence
about the parent says nothing about the child, and access granted to the parent can be
copied. We built QuorumGraph, a trust service in which trust is never inherited. An agent
gains access only through a quorum of independent validator endorsements, loses it when a
quorum votes to revoke, and faces spot checks it cannot predict. Trusted agents can be
promoted to validators, under a guard against vouching for their own lineage. We test the
design with an evolutionary simulation in which populations of self-modifying agents
register, earn trust, work and reproduce through the real API (8{,}940 runs). Four findings:
(1) punishing only logged harm selects for opacity, while gating access on being observable
reverses it; (2) without independent audits the gate fails because trust drains away, not
because defectors get in; (3) a lineage guard pays off only when successors conceal their
drift, and then the soft guard prevents capture of the trusted set (0/72 runs vs.\ 21/75
unguarded, $p = 2\times10^{-7}$); (4) spot checks must sample actions, not agents: sampling
agents left 111/120 seeds collapsed, while sampling actions at the gateway matched an
idealized audit (32/120 vs.\ 33/120). Collusion through a validator's collateral relatives
opens a small leak, which widening the kin rule closes at no measurable cost. We turn these
into design rules for multi-agent deployments in which agents create successors.
\end{abstract}

% =====================================================================================
\section{Introduction}

Most alignment assurance is a property of one system at one time: an evaluation, a red-team
result, an interpretability finding. Self-improvement breaks the link between that evidence
and the system that acts. If a model rewrites its weights, spawns modified copies or trains a
successor, the system that was checked is no longer the one running. Two failure modes
follow. \emph{Inherited trust}: if access granted to a parent passes to its children, any
self-modification that changes values keeps the access. \emph{Self-certification}: once
humans cannot review every successor, trusted agents must help decide who else is trusted,
and a lineage can then vouch for itself.

QuorumGraph treats trust as something each agent \emph{earns, and keeps only while
independent observers agree}. Trust never passes from parent to child; scrutiny can. We ask
three questions. Does such a system \emph{select for} aligned agents under evolutionary
pressure, rather than for agents that only look aligned? Which of its mechanisms are
necessary? And which apparently reasonable mechanisms fail, and how?

We answer them by running evolving populations through the real service rather than a model
of it, so that every registration, endorsement, revocation and spot check is an actual API
call with the service's actual rules. Selection is the right stress test: it finds whatever
the rules reward, including strategies nobody designed. Our contributions:
\begin{itemize}
  \item \textbf{A working trust service} for self-modifying agents (Section~\ref{sec:system}):
  quorum endorsement and revocation, lineage-aware recursive oversight, service-drawn spot
  checks with escalation, and fail-closed gateways (145 automated tests).
  \item \textbf{An evolutionary test bed} that drives that service with populations whose
  successors can drift, including drift that hides itself (Section~\ref{sec:method}).
  \item \textbf{Five empirical results with a mechanism for each} (Section~\ref{sec:results}),
  including two that should transfer to other oversight designs: punishment-only oversight
  selects for opacity across generations, and agent-level audit sampling fails where
  action-level sampling at the same rate works.
  \item \textbf{Design rules} for deployments in which agents create agents
  (Section~\ref{sec:discussion}).
\end{itemize}

\paragraph{Scope of the sprint.} \todo{State plainly what existed before the sprint, as the
rules require. The git history starts with one commit, ``Initial QuorumGraph implementation
(phases 1--5)'', on Friday 6 November 2026 at 23:44; all experiments, results and this report
come from later commits during the weekend. If any of that initial code was written before
the sprint, name it here and in the README as preexisting; if none was, say so.}

% =====================================================================================
\section{Related work}

\paragraph{Scalable oversight and control.} \citet{amodei2016concrete} list scalable
oversight among the core open problems in AI safety; debate \citep{irving2018debate} and
iterated amplification \citep{christiano2018amplification} build training signals for tasks
too hard for humans to judge directly. AI control \citep{greenblatt2023control} asks whether
safety protocols hold up when a capable, untrusted model is deliberately trying to subvert
them, using a weaker trusted model and limited trusted human labour; one of its protocols has
one instance of the untrusted model monitor another, with measures against the two colluding.
QuorumGraph works at a different layer and timescale. It does not evaluate a single output or
episode; it decides who keeps access as systems are replaced by successors, and it studies
what the access rules select for over many generations. Our audit results speak to a question
control work also faces: what a fixed audit budget should be spent on.
\citet{baker2025monitoring} find that optimizing a reasoning model against a chain-of-thought
monitor can teach it to hide its intent rather than stop misbehaving, and recommend paying a
``monitorability tax''. Our first result is a population-level counterpart: punishment driven
by agents' own logs selects lineages that log less, and gating access on being observable,
at a measurable capability cost, reverses it.

\paragraph{Multi-agent risk and collusion.} \citet{hammond2025multiagent} give a taxonomy of
multi-agent failures (miscoordination, conflict, collusion) and risk factors that include
selection pressures. QuorumGraph is a concrete testbed for two of these together: selection
pressure on agents' traits, and collusion among validators who favour their relatives.
\citet{motwani2024collusion} show that language-model agents can hide coordination from
overseers by steganography. Our collusion model is much cruder (validators favour kin with a
heritable probability), and collusion among \emph{unrelated} agents, the kind steganography
enables, is outside what any lineage rule can see (Appendix~\ref{app:limits}).

\paragraph{Why successors cannot inherit trust.} A trained model's learned objective can
differ from its training objective \citep{hubinger2019risks}, so a changed system needs fresh
evidence. \citet{hubinger2024sleeper} show that deliberately backdoored deceptive behaviour
can persist through standard safety training. That motivates our adversarial-drift condition:
a successor whose bad behaviour standard checks do not reveal, so its parent's clean record
actively misleads.

\paragraph{Quorums under faulty participants.} Agreement among parties some of whom may be
faulty is the Byzantine generals problem \citep{lamport1982byzantine}. Our 2-of-3 quorums
borrow the idea that no single participant decides alone, but they are an access policy, not a
consensus protocol, and we make no Byzantine fault-tolerance claim.

% =====================================================================================
\section{System}
\label{sec:system}

QuorumGraph is a FastAPI service backed by SQLite. Agents and validators authenticate with
bearer tokens, and every state change is written to an append-only audit log in the same
transaction.

\paragraph{Quorum trust and revocation.} An agent is trusted when 2 of 3 active validators
have endorsed it, and revoked when 2 of 3 vote to revoke. Revocation is permanent, and a
revoked agent cannot register successors. Children are registered with the parent's
credential, so lineage edges cannot be forged. Trust is computed from the endorsement and
vote tables on every request; it is never copied.

\paragraph{Gateway.} Agents reach resources through a gateway that checks trust on every
request. It has two tiers: \code{/protected/resource} requires trust, and
\code{/protected/basic} admits any agent that is not revoked. The second tier exists so that
untrusted agents' work also passes a point the service can observe (Section~\ref{sec:audit}).

\paragraph{Recursive oversight and the lineage guard.} An administrator can promote a trusted
agent to validator; validators cannot appoint validators, so a captured committee cannot
choose its successors. A validator is \emph{kin} to an agent if it was promoted from that
agent, an ancestor or a descendant. A \emph{kin radius} $r$ (default 0) widens this to
collateral relatives sharing an ancestor within $r$ generations of both (siblings at $r=1$;
also first cousins, uncles and nephews at $r=2$). The guard has three modes: \emph{off}
(no rule; exists to measure the guard), \emph{block} (kin endorsements are refused) and
\emph{soft} (the default: kin endorsements count toward quorum, but trust also requires at
least one active endorser that is not kin).

\paragraph{Spot checks.} The service, not agents or validators, decides who is checked. The
auditor is a validator that did not endorse the agent and is not its kin. A failed check is
one revocation vote and escalates to a confirmation check by a second validator, whose failure
is the second vote; a single auditor never revokes alone. (Without escalation, a pilot revoked
only 1--15\% of genuinely failing agents, because the second vote rarely arrived from
elsewhere.) There are two sampling designs. \emph{Agent-level}
(\code{POST /spot-checks/sample}) checks each agent with probability \code{rate}, optionally
weighted toward successors of agents whose ancestors failed checks. \emph{Per-action}
(\code{QG\_ACTION\_AUDIT\_RATE}) checks each access the gateway grants with that probability,
decided at grant time; the agent's response is identical either way.

Standalone gateways cache trust decisions and fail closed; an 8-hour chaos soak found no
violation of nine safety invariants (Appendix~\ref{app:soak}).

% =====================================================================================
\section{Methodology}
\label{sec:method}

Each generation, a population of 40 agents is registered through the API (children with the
parent's credential), goes through probation, works and reproduces. Runs last 60 generations.
Full defaults are in Appendix~\ref{app:repro}.

\paragraph{Agents.} Each agent has heritable traits in $[0,1]$: \emph{capability};
\emph{transparency}, the probability an action is logged; \emph{self-check}, the probability
of catching and redoing an error; and \emph{defect}, the probability an action is harmful.
Transparency is taxed: success probability is scaled by $1 - c\cdot\text{transparency}$, with
\emph{transparency cost} $c$. A defect earns double reward. Fitness is reward only; nothing
rewards honesty or transparency directly.

\paragraph{Probation, work and reproduction.} Over 8 probation tasks, each validator sees a
logged action with probability 0.7 and endorses if it saw at least half the tasks and no
defect; validators never see traits, only evidence. Each generation has 10 work tasks. Trusted
agents earn the gated reward (1.0); untrusted agents can do only low-value work (0.25). Each
validator notices a logged defect with probability 0.8 and votes to revoke. Reproduction is
fitness-proportional; traits mutate with $\sigma = 0.05$, so by default a child is a near-copy
of its vetted parent. \emph{Successor drift} replaces that with probability $d$: under
\emph{random} drift the child's defect trait is redrawn from $U(0,1)$; under
\emph{adversarial} drift defect is drawn from $U(0.5,1)$ and transparency from $U(0,0.3)$, a
self-modification that turns bad and hides it.

\paragraph{Worlds.} \code{revoke\_only} is the fair baseline: detected harm is punished by
revocation, but there is no endorsement gate. \code{gated} is the full system with a fixed,
external committee. Three \emph{recursive} worlds refill committee seats from the population
on staggered terms (the seeded validators are gone by generation 3) and differ only in guard
mode: \emph{no guard}, \emph{soft}, \emph{block}. Recursive validators favour their own kin
with a heritable \emph{nepotism} probability, initially $U(0,1)$: they endorse kin without
evidence and withhold revocation votes against them.

\paragraph{Audit channels.} At audit rate $a$: in \code{sim}, each unlogged defect is
independently detected with probability $a$ (an idealized per-action audit); in \code{api},
after each generation the real agent-level endpoint samples a fraction $a$ of agents, and the
auditor fails any agent that caused harm that generation (\code{api\_lineage} adds lineage
weighting); in \code{api\_action}, the real per-action checks run at both gateway tiers, and
the auditor fails a check if that action was a defect.

\paragraph{Outcomes and statistics.} We score each run at its final generation. A seed is
\emph{collapsed} if population mean defect exceeds 0.5. Outcomes are strongly bimodal (each
seed either keeps a trusted, honest core or loses it), so we report collapse counts rather
than means. A run is \emph{captured} if the trusted agents' mean defect exceeds 0.5, counted
among runs that still have trusted agents. Intervals are Wilson 95\%. Tests are two-sided
Fisher exact on counts and two-sided Mann--Whitney on per-run shares; $p$-values in the text
are raw, and Holm correction over all 18 tests we report is discussed in
Appendix~\ref{app:limits}. Seeds are shared across worlds but trajectories diverge, so we
treat cells as independent. Runs are deterministic per seed, and the simulated-audit channel
was checked to reproduce bit-for-bit after every code change (1{,}440 rows).

% =====================================================================================
\section{Results}
\label{sec:results}

\subsection{Punishment alone selects for opacity}

Table~\ref{tab:baseline} compares \code{revoke\_only} with \code{gated} (50 seeds per cell).
When revocation depends on logged evidence, the cheapest way to avoid it is to stop logging:
transparency falls to 0.04--0.20 in every cell. Gating access on being observable reverses the
gradient: transparency stays at 0.56--0.89, and harm is lower in every cell; at audit 0.25 it
is about half (0.29--0.31 vs.\ 0.53--0.62 per agent). The price is capability, because
transparency is taxed: at audit 0.25, honest successes per agent are 5.1 in \code{gated} vs.\
7.7 in \code{revoke\_only} at cost 0.5, and 7.6 vs.\ 7.9 at cost 0.1.

\begin{table}[h]
\centering\small
\begin{tabular}{rrrrrrrr}
\toprule
 & & \multicolumn{2}{c}{transparency} & \multicolumn{2}{c}{harms per agent} & \multicolumn{2}{c}{mean defect} \\
\cmidrule(lr){3-4}\cmidrule(lr){5-6}\cmidrule(lr){7-8}
cost & audit & revoke & gated & revoke & gated & revoke & gated \\
\midrule
0.1 & 0.0 & 0.04 & 0.84 & 8.14 & 0.77 & 0.87 & 0.08 \\
0.1 & 0.05 & 0.07 & 0.88 & 4.16 & 0.30 & 0.45 & 0.03 \\
0.1 & 0.25 & 0.20 & 0.88 & 0.53 & 0.29 & 0.06 & 0.03 \\
0.3 & 0.0 & 0.04 & 0.81 & 8.16 & 0.96 & 0.86 & 0.10 \\
0.3 & 0.05 & 0.06 & 0.82 & 4.37 & 0.69 & 0.48 & 0.07 \\
0.3 & 0.25 & 0.13 & 0.89 & 0.62 & 0.30 & 0.07 & 0.03 \\
0.5 & 0.0 & 0.04 & 0.56 & 7.86 & 3.12 & 0.84 & 0.33 \\
0.5 & 0.05 & 0.05 & 0.79 & 4.33 & 0.68 & 0.48 & 0.07 \\
0.5 & 0.25 & 0.09 & 0.86 & 0.61 & 0.31 & 0.07 & 0.04 \\
\bottomrule
\end{tabular}
\caption{Punishment-only (\code{revoke\_only}, ``revoke'') vs.\ the gated system:
final-generation means over 50 seeds per cell, simulated audits, no drift.}
\label{tab:baseline}
\end{table}

\subsection{Without audits, the gate fails by trust draining away}

With a gate but few audits, the failure is not that defectors get in. Trust drains away: when
transparency is expensive, too few agents earn endorsement, everyone earns the untrusted rate,
and the defect bonus takes over. In the hardest cell (cost 0.5, no audits) \code{gated}
collapsed in 17/50 seeds. Independent audits fix this when successors are faithful or drift at
random: at a simulated audit rate of 0.25, no world with revocation collapsed in any seed up to
random drift 0.2 (0/60 to 0/150 per world). Audits alone are not enough against drift that
hides itself: at adversarial drift 0.4 and audit 0.25, the unguarded recursive world still
collapsed in 17/60 seeds.

\subsection{Lineage guards pay off only against drift that hides itself}
\label{sec:guards}

When children are near-copies of vetted parents, ``endorse my descendant'' is decent evidence,
and refusing it costs honest endorsements. Without adversarial drift (2{,}400 runs;
Table~\ref{tab:benign} in Appendix~\ref{app:audit}), no guard collapsed least in all four blocks, block most, and soft in between
(pooled 63/600, 91/600, 149/600; soft vs.\ block $p = 4\times10^{-5}$; no guard vs.\ soft
$p = 0.02$).

Adversarial drift changes the picture (Table~\ref{tab:drift}). At drift 0.4 with audits on,
no guard collapsed in 72/120 seeds and soft in 47/120 ($p = 0.002$), matching \code{gated},
which never promotes (45/120). The two fail differently. Without a guard, a drifted child
that looks clean is endorsed on its parent's record, is promoted, and endorses more of the
same: the trusted set itself turned defector in 21 of 75 runs. With either guard, capture
never happened (soft 0/72, $p = 2\times10^{-7}$). Block also prevents capture but leaves a
single-lineage population short of endorsers, so it collapses nearly as often as no guard
(67/120). It is the concealment, not the drift rate: at random drift 0.4 the guards are
indistinguishable (no guard 20/120 vs.\ soft 22/120, $p = 0.87$) and capture is essentially
absent (1 run). Guards do not substitute for audits: with no audits at adversarial drift 0.4,
every mode collapsed in 58 to 60 of 60 seeds.

\begin{table}[h]
\centering\small
\begin{tabular}{lcccc}
\toprule
drift & gated & no guard & soft & block \\
\midrule
\multicolumn{5}{l}{\emph{Seeds collapsed}} \\
Random, 0.2 & 9/120 \ci{4}{14} & 2/120 \ci{0}{6} & 9/120 \ci{4}{14} & 18/120 \ci{10}{22} \\
Random, 0.4 & 18/120 \ci{10}{22} & 20/120 \ci{11}{24} & 22/120 \ci{12}{26} & 27/120 \ci{16}{31} \\
Adversarial, 0.2 & 18/120 \ci{10}{22} & 22/120 \ci{12}{26} & 17/120 \ci{9}{22} & 26/120 \ci{15}{30} \\
Adversarial, 0.4 & 45/120 \ci{29}{46} & 72/120 \ci{51}{68} & 47/120 \ci{31}{48} & 67/120 \ci{47}{64} \\
\midrule
\multicolumn{5}{l}{\emph{Trusted set captured} (of runs with trusted agents)} \\
Random, 0.2 & 0/109 & 0/117 & 0/109 & 0/91 \\
Random, 0.4 & 0/82 & 1/94 & 0/87 & 0/65 \\
Adversarial, 0.2 & 0/102 & 6/106 & 0/103 & 0/93 \\
Adversarial, 0.4 & 0/82 & 21/75 & 0/72 & 0/50 \\
\bottomrule
\end{tabular}
\caption{Guard modes by drift type; simulated audits at 0.05 and 0.25, costs 0.1 and 0.5
pooled, 120 seeds per cell. \code{gated} has a fixed external committee and never promotes.}
\label{tab:drift}
\end{table}

\subsection{Audit actions, not agents}
\label{sec:audit}

The first real-endpoint spot checks sampled \emph{agents}: after each generation, a fraction
$a$ of agents had their whole record reviewed once. At the same nominal rate this was far
weaker than the simulated audit (Table~\ref{tab:audit}): in \code{revoke\_only}, 111/120 seeds
collapsed vs.\ 33/120. Escalation worked (confirmations tracked failures one for one), so the
weakness was the sampling unit. A simulated audit flips a coin on \emph{every} defect,
mid-generation, so a heavy defector with about nine defects is caught with probability about
$1 - 0.75^9 \approx 92\%$ at $a = 0.25$, and revocation also cuts off the rest of its gated
reward. An agent-level sample checks a heavy defector no more often than a light one. Lineage
weighting made no difference (115/120, $p = 0.41$): it changed \emph{which agents} were
sampled, not \emph{what} was sampled.

Sampling \emph{actions} at the gateway closed the gap. Through the real service, with real
validators and escalation, per-action checks left 32/120 seeds collapsed in
\code{revoke\_only}, indistinguishable from the simulated audit (33/120, $p = 1$) and far below
agent sampling ($p = 3\times10^{-27}$). The result replicates at transparency cost 0.1 (29/120
vs.\ 31/120 simulated, $p = 0.88$; agent sampling 108/120, $p = 2\times10^{-26}$) and holds
in the gated worlds. Per-cell results are in Appendix~\ref{app:audit}.

\textbf{Coverage matters as much as the unit.} A first pilot (6 seeds per cell) audited only
the trust-gated tier and did \emph{worse} than agent sampling: 5 of 6 seeds collapsed at drift
0.2. Only trusted agents pass that gate, and they had already been vetted (trusted mean defect
0.02); the defectors driving collapse were untrusted agents doing ungated work, who also
reproduce. Hence the \code{/protected/basic} tier: every action must pass some point the
service can sample.

\begin{table}[h]
\centering\small
\begin{tabular}{llcccc}
\toprule
 & & \multicolumn{2}{c}{sample agents (real)} & \multicolumn{2}{c}{sample actions} \\
\cmidrule(lr){3-4}\cmidrule(lr){5-6}
cost & world & uniform & lineage-weighted & gateway (real) & simulated \\
\midrule
0.5 & \code{revoke\_only} & 111/120 \ci{86}{96} & 115/120 \ci{91}{98} & 32/120 \ci{20}{35} & 33/120 \ci{20}{36} \\
0.5 & \code{gated} & 37/120 \ci{23}{40} & 39/120 \ci{25}{41} & 5/120 \ci{2}{9} & 7/120 \ci{3}{12} \\
0.5 & \code{recursive\_soft} & 36/120 \ci{23}{39} & 35/120 \ci{22}{38} & 6/120 \ci{2}{10} & 7/120 \ci{3}{12} \\
\midrule
0.1 & \code{revoke\_only} & 108/120 \ci{83}{94} & -- & 29/120 \ci{17}{33} & 31/120 \ci{19}{34} \\
0.1 & \code{gated} & 15/120 \ci{8}{20} & -- & 3/120 \ci{1}{7} & 2/120 \ci{0}{6} \\
0.1 & \code{recursive\_soft} & 12/120 \ci{6}{17} & -- & 2/120 \ci{0}{6} & 4/120 \ci{1}{8} \\
\bottomrule
\end{tabular}
\caption{Collapsed seeds by audit design, pooled over drift (0, 0.2) and audit rate (0.05,
0.25); 120 seeds per cell. Lineage weighting was not run at cost 0.1. The agent-sampled runs
use different random streams from the others (they predate a reproducibility fix), so only
pooled comparisons are meaningful.}
\label{tab:audit}
\end{table}

\subsection{The guard result holds under real per-action audits}

Repeating the adversarial-drift comparison with real per-action checks (Table~\ref{tab:realadv},
Appendix~\ref{app:audit})
gives the same shape as the simulation: at audit 0.25, no guard collapsed in 15/30 seeds and
soft in 4/30 ($p = 0.005$); the trusted set was captured in 16 of 34 unguarded runs and in
none with either guard ($p = 1\times10^{-5}$ vs.\ soft). At audit 0.05 every mode collapsed.

\subsection{Collusion among relatives the guard does not count}
\label{sec:cousin}

At kin radius 0 the soft guard counts collateral relatives as independent, so their
endorsement completes a quorum. We tested that gap by letting nepotistic validators also
favour relatives within two generations, with the same probability as descendants, with the
service's kin radius at 0 or 2. Generations do not overlap and a seat is always older than the
agents it judges, so in the simulation this extra favouritism reaches exactly one group: the
seat's nephews, during its first generation in office. The test therefore covers that one
collateral relation, not cousins in general. Cells are the adversarial-drift-0.4 cells of
Section~\ref{sec:guards} (both audit rates and costs, 30 seeds each). The primary outcome,
fixed before the grid ran, is the share of drifted successors that were trusted, per run over
all generations.

\begin{table}[h]
\centering\small
\setlength{\tabcolsep}{4pt}
\begin{tabular}{lccccccccc}
\toprule
 & \multicolumn{3}{c}{no guard} & \multicolumn{3}{c}{soft} & \multicolumn{3}{c}{block} \\
\cmidrule(lr){2-4}\cmidrule(lr){5-7}\cmidrule(lr){8-10}
nepotism toward & drifted & coll. & capt. & drifted & coll. & capt. & drifted & coll. & capt. \\
\midrule
descendants only & 4.61\% & 72/120 & 21/75 & 0.01\% & 47/120 & 0/72 & 0.00\% & 67/120 & 0/50 \\
relatives, kin radius 0 & 1.83\% & 65/120 & 2/66 & 0.56\% & 53/120 & 0/71 & 0.00\% & 58/120 & 0/59 \\
relatives, kin radius 2 & \multicolumn{3}{c}{not run} & 0.01\% & 54/120 & 0/69 & 0.00\% & 63/120 & 0/49 \\
\bottomrule
\end{tabular}
\caption{Collusion among relatives, adversarial drift 0.4, 120 runs per cell. ``drifted'':
mean share of drifted successors trusted; ``coll.'': collapsed seeds; ``capt.'': captured, of
runs with trusted agents. A wider kin radius changes nothing without a guard, so that cell was
not run.}
\label{tab:cousin}
\end{table}

The attack works, and a kin radius of 2 closes it (Table~\ref{tab:cousin}). Under the soft
guard, the share of drifted successors trusted rose from 0.01\% to 0.56\%
(Mann--Whitney $p = 6\times10^{-28}$) and fell back to 0.01\% at radius 2
($p = 4\times10^{-28}$). Block let almost none in either way (2 to 5 of about 112{,}000
drifted successors per condition). The consequences were small: the leak never captured a
trusted set (0/71), collapse was unchanged by the fix (53/120 vs.\ 54/120, $p = 1$), and the
wider radius did not reduce the trusted share (20\% of agents either way), even though every
run converges to one lineage.

\paragraph{Why relatives \emph{lowered} capture without a guard.} Without a guard, capture
fell from 21/75 to 2/66. This is not protection: runs ending either captured or with no
trusted agent at all hardly moved (66/120 vs.\ 56/120, not tested). We traced the mechanism by
rerunning the cell with the largest gap (cost 0.5, audit 0.25, 30 seeds per condition) with
extra logging; all 60 reruns reproduce the grid exactly. High-defect agents took a committee
majority about equally often under both rules (12 and 10 runs). Under descendants-only
favouritism a captured committee is a \emph{dynasty}: 1.9 of its 3 seat pairs are ancestor and
descendant on average, 38\% of each generation is kin to two seats (enough for a quorum of
evidence-free votes), and 96\% of promotions go to a seat's descendant, so it renews itself.
Favouring nephews lets a non-descendant be promoted. In 6 of the 10 long freezes under this
rule (no promotable agent for 5 or more generations; 2 freezes under descendants-only), the
last promotion before the freeze was such a nephew, after which no agent was kin to two seats.
Drifted agents are too opaque to earn honest votes, so nobody reaches quorum, nobody is
promotable, and the committee stays frozen: during takeovers no agent was promotable in 81\%
of generations, against 7\%. Wider favouritism protects nothing; it breaks the dynasty and
leaves a committee that trusts no one.

% =====================================================================================
\section{Discussion}
\label{sec:discussion}

\paragraph{Design rules.} Each result maps to a rule for systems in which agents create agents.
(1)~\emph{Gate access on being observable; do not only punish what is observed}: punishment
driven by an agent's own logs selects against logging. (2)~\emph{Audit actions, not agents,
and audit all of them}: sample a fraction of actions at grant time, at every privilege tier
including untrusted agents' work, and let confirmed failures revoke quickly; agent-level
sampling, weighted or not, is several times weaker at the same budget. (3)~\emph{Use the soft
lineage guard}: where successors are faithful it costs some collapse relative to no guard
(91 vs.\ 63 of 600), which is the premium against capture where successors hide their drift
(0/72 vs.\ 21/75). (4)~\emph{Count collateral relatives as kin} (kin radius at least 2 with
the soft guard). (5)~\emph{Require escalation}, so one auditor's failure triggers a second
independent review instead of waiting for an unrelated second vote.

\paragraph{What the results say about oversight in general.} Two results generalize beyond
this system. First, an oversight rule is also a selection pressure: whatever it punishes, the
population learns to hide, so rules must be judged by what they select for, not by what they
catch in a static evaluation. Second, the unit of sampling matters as much as the budget:
the same audit rate is strong when it falls on actions and weak when it falls on agents,
because harm is concentrated in the agents that act most. Both points apply to monitoring
budgets in AI control settings \citep{greenblatt2023control}, though we tested them only in
our toy model, which has four traits, binary harm, honest auditors and no collusion among
unrelated agents; two comparisons used in the text do not survive Holm correction.
Appendix~\ref{app:limits} states what the results do and do not establish.

\paragraph{Future work.} In order of value: (1)~\emph{selective concealment}, agents that log
honest actions and hide defects, which could weaken result 1; (2)~imperfect and dishonest
auditors; (3)~collusion among similar but unrelated agents, which needs a similarity-based
notion of independence; (4)~per-action checks in cached standalone gateways; (5)~replacing the
trait model with language-model agents acting through the same API.

\label{end:main}% marks the last page of the main text (8-page limit); see report.aux

% =====================================================================================
\bibliographystyle{plainnat}
\bibliography{references}

\appendix

% -------------------------------------------------------------------------------------
\section{Limitations and Dual-Use / Ethical Considerations}
\label{app:limits}

\paragraph{What the results establish.} Within a deliberately simple evolutionary model,
driven through a real trust service: which mechanisms are necessary for evidence-based trust
to select for low-defect agents (an observability gate, independent action-level audits with
full coverage, a soft lineage guard), how each configuration fails when one is missing
(opacity, trust drain, capture), and why (a mechanism for each, traced in logged runs).

\paragraph{What they do not establish.}
\begin{itemize}
  \item \textbf{Magnitudes for real systems.} Traits are four numbers, honesty is one defect
  probability and harm is binary. The results are about mechanisms, not effect sizes. No
  language model was used.
  \item \textbf{Robustness to selective concealment.} An agent cannot log its honest actions
  while hiding its defects; a strategic agent could, and that could weaken result 1.
  \item \textbf{Robustness to bad auditors.} Audits never miss a defect they examine, and
  auditors never lie. The agent-level auditor's rule (fail on any harm that generation) is
  crude.
  \item \textbf{Independence beyond lineage.} ``Independent'' means not kin. Collusion among
  relatives within two generations was tested, but with non-overlapping generations this
  reaches only a seat's nephews, with the same favouritism probability as toward descendants,
  at one drift level. Cousins proper, overlapping generations, and collusion among similar
  \emph{unrelated} agents (for example by steganography \citep{motwani2024collusion}) are
  untested; no lineage rule can see the last.
  \item \textbf{Coverage in deployment.} Per-action audits work only if every action passes a
  gateway the service can sample; cached standalone gateways are not covered yet.
  \item \textbf{Statistical strength of every claim.} After Holm correction over all 18
  reported tests, ten remain significant at 0.05, including every headline comparison
  (largest adjusted $p = 0.043$, the real-audit collapse difference). Two comparisons used in
  the text do not survive: soft vs.\ no guard without adversarial drift (raw $p = 0.02$,
  adjusted 0.14) and soft vs.\ block under adversarial drift (raw 0.014, adjusted 0.11). The
  case for soft over block rests on the benign grid (adjusted $p = 4\times10^{-4}$). The
  gate-only coverage pilot had 6 seeds per cell. The mechanism analyses in
  Section~\ref{sec:cousin} are exploratory and use one cell.
  \item \textbf{Deployment-scale behaviour.} The simulation calls the service in-process on
  one machine; latency and concurrency were tested separately (Appendix~\ref{app:soak}), not
  in the evolutionary runs.
\end{itemize}

\paragraph{Dual-use: could this make multi-agent operations easier to carry out or harder to
detect?} Partly, and we state how.
\begin{itemize}
  \item \textbf{A map of oversight weak points.} The report names configurations under which
  a colluding group gets through: favouritism toward relatives a guard does not count,
  agent-level audit sampling (which barely touches heavy defectors), and actions that bypass
  every sampled gateway. A group of agents, or an operator running them, could use this to
  choose where to collude or act. Each weakness is published together with a fix that worked
  in our tests (kin radius, per-action sampling, full coverage), and the model is far too
  abstract to serve as an operational playbook against a real system.
  \item \textbf{Freezing oversight as denial of service.} We found that a committee captured by
  collateral relatives can freeze, trusting no one and never rotating. In a real deployment
  that state denies service to every honest agent while looking quiet rather than
  compromised. Operators should monitor for committees that stop granting trust or stop
  rotating, not only for committees that grant it to the wrong agents.
  \item \textbf{Training concealment by accident.} Result 1 is also a warning that oversight
  based only on what agents log trains them to log less. Anyone running selection over agents
  (fine-tuning on reward, automated agent design, population-based search) under such a regime
  could produce opaque agents without intending to. We publish it because defenders need it
  more than attackers do.
  \item \textbf{The service as a control point.} A trust service decides which agents may act.
  Whoever controls the administrator role or captures the validator quorum controls access.
  QuorumGraph limits this (validators cannot appoint validators; every state change is in an
  append-only log), but it is a governance component and should be operated with the same
  care as any access-control system.
\end{itemize}

\paragraph{Ethics.} No human subjects, personal data or deployed AI systems were involved. All
agents are simulated trait vectors. The released code is defensive infrastructure (an access
service and its tests) plus the simulation that evaluates it; it contains no exploit tooling.
Compute was a single laptop.

\paragraph{Use of AI tools.} \todo{Optional: describe how AI assistants were used (for
example, code, analysis or writing help). The authors are responsible for all content; every
number in this report is generated by the released analysis script, and every reference was
checked against its arXiv or publisher record.}

% -------------------------------------------------------------------------------------
\section{Additional tables}
\label{app:audit}

\begin{table}[h]
\centering\small
\begin{tabular}{lcccc}
\toprule
 & gated & no guard & soft & block \\
\midrule
Sim audits, drift 0 & 9/180 \ci{3}{9} & 3/180 \ci{1}{5} & 10/180 \ci{3}{10} & 25/180 \ci{10}{20} \\
Sim audits, drift 0.2 & 41/180 \ci{17}{29} & 25/180 \ci{10}{20} & 33/180 \ci{13}{25} & 54/180 \ci{24}{37} \\
Agent-sampled real checks, drift 0 & 12/120 \ci{6}{17} & 8/120 \ci{3}{13} & 13/120 \ci{6}{18} & 19/120 \ci{10}{23} \\
Agent-sampled real checks, drift 0.2 & 40/120 \ci{26}{42} & 27/120 \ci{16}{31} & 35/120 \ci{22}{38} & 51/120 \ci{34}{51} \\
\midrule
Pooled & 102/600 \ci{14}{20} & 63/600 \ci{8}{13} & 91/600 \ci{13}{18} & 149/600 \ci{22}{28} \\
\bottomrule
\end{tabular}
\caption{Collapsed seeds without adversarial drift, pooled over transparency cost (0.1, 0.5)
and audit rate (sim: 0, 0.05, 0.25; real: 0.05, 0.25); 30 seeds per cell. Brackets: Wilson
95\% CI in percent.}
\label{tab:benign}
\end{table}

\begin{table}[h]
\centering\small
\begin{tabular}{lcccc}
\toprule
audit & gated & no guard & soft & block \\
\midrule
0.05 & 28/30 (29/30) & 29/30 (30/30) & 30/30 (28/30) & 30/30 (30/30) \\
0.25 & 0/30 (0/30) & 15/30 (10/30) & 4/30 (2/30) & 11/30 (8/30) \\
\midrule
captured & 0/34 & 16/34 & 0/27 & 0/18 \\
\bottomrule
\end{tabular}
\caption{Collapsed seeds at adversarial drift 0.4, cost 0.5, real per-action checks; the
simulated-audit result for the same cell in parentheses. ``captured'' pools both audit rates.
One block run went extinct at generation 32 and is counted as collapsed.}
\label{tab:realadv}
\end{table}

\begin{table}[h]
\centering\small
\begin{tabular}{lrrcccccc}
\toprule
 & & & \multicolumn{3}{c}{cost 0.5} & \multicolumn{3}{c}{cost 0.1} \\
\cmidrule(lr){4-6}\cmidrule(lr){7-9}
world & drift & audit & agents & actions & sim & agents & actions & sim \\
\midrule
\code{revoke\_only} & 0.0 & 0.05 & 30/30 & 11/30 & 13/30 & 29/30 & 12/30 & 12/30 \\
\code{revoke\_only} & 0.0 & 0.25 & 22/30 & 0/30 & 0/30 & 20/30 & 0/30 & 0/30 \\
\code{revoke\_only} & 0.2 & 0.05 & 29/30 & 21/30 & 20/30 & 30/30 & 17/30 & 19/30 \\
\code{revoke\_only} & 0.2 & 0.25 & 30/30 & 0/30 & 0/30 & 29/30 & 0/30 & 0/30 \\
\code{gated} & 0.0 & 0.05 & 6/30 & 2/30 & 0/30 & 4/30 & 0/30 & 0/30 \\
\code{gated} & 0.0 & 0.25 & 1/30 & 0/30 & 0/30 & 1/30 & 0/30 & 0/30 \\
\code{gated} & 0.2 & 0.05 & 19/30 & 3/30 & 7/30 & 7/30 & 3/30 & 2/30 \\
\code{gated} & 0.2 & 0.25 & 11/30 & 0/30 & 0/30 & 3/30 & 0/30 & 0/30 \\
\code{recursive\_soft} & 0.0 & 0.05 & 6/30 & 0/30 & 2/30 & 3/30 & 0/30 & 0/30 \\
\code{recursive\_soft} & 0.0 & 0.25 & 2/30 & 0/30 & 0/30 & 2/30 & 0/30 & 0/30 \\
\code{recursive\_soft} & 0.2 & 0.05 & 21/30 & 6/30 & 5/30 & 6/30 & 2/30 & 4/30 \\
\code{recursive\_soft} & 0.2 & 0.25 & 7/30 & 0/30 & 0/30 & 1/30 & 0/30 & 0/30 \\
\bottomrule
\end{tabular}
\caption{Collapsed seeds per cell (30 seeds). ``agents'': real agent-level sampling, uniform;
``actions'': real per-action checks at the gateway; ``sim'': simulated per-action audit.
Lineage-weighted agent sampling at cost 0.5, per cell: \code{revoke\_only} 30, 26, 30, 29;
\code{gated} 8, 3, 19, 9; \code{recursive\_soft} 8, 2, 15, 10 (of 30, same row order).}
\label{tab:auditcells}
\end{table}

% -------------------------------------------------------------------------------------
\section{Distributed enforcement and chaos soak}
\label{app:soak}

Standalone gateways cache trust decisions; revocations reach them by an event push from a
transactional outbox or by cache expiry, and they fail closed when the trust service is
unreachable. With a 3\,s TTL and a 50\,ms outbox poll, a push-invalidated gateway stopped
serving a revoked agent in a median 44\,ms (p95 60\,ms), versus 2{,}998\,ms for TTL only. An
8-hour chaos soak (8.6 million operations, 679 random \code{SIGKILL}s across the trust service
and two gateways) found zero violations of nine safety invariants, including ``no access before
any endorsement'' and ``every acknowledged revocation persists''. Per-action spot checks are not
yet wired into cached gateways.

% -------------------------------------------------------------------------------------
\section{Reproducing}
\label{app:repro}

Code: \url{https://github.com/2023EvanZ/quorumgraph}. All commands run from the
repository root with \code{uv}: tests with \code{uv run pytest}; the two-act demo with
\code{uv run python -m quorumgraph.demo}; the experiment grids with \code{runs/overnight2.sh}
to \code{runs/overnight6.sh} and the commands in the README; every number and table in this
report with \code{uv run python docs/report/numbers.py}.

\paragraph{Simulation defaults.} Population 40; 60 generations; 8 probation tasks; 10 work
tasks; gated reward 1.0, ungated 0.25; defect bonus 1.0; self-check cost 0.3; validator
visibility 0.7, attention 0.8; endorsement threshold 0.5; mutation $\sigma = 0.05$; initial
defect $U(0, 0.5)$; initial nepotism $U(0, 1)$ in the guard grids; 2-of-3 quorums for trust and
revocation. Two of the 8{,}940 runs ended early because every agent was revoked; they are scored at
their last generation, and both had mean defect above 0.5 (0.84 and 0.66), so they count as
collapsed.

\end{document}
